% Einheit 1 von 4 – Punktprobe und Seitenlage (bank/lagebeziehungen/e1.jsonl, zusammenbau v0.1)
%% TODO Zeitmarke Einheit 1: Marken-Zeile nicht lesbar
%% TODO Zweigzeile Teil 1: was hier gelernt wird (Satz in Schülersprache aus dem Einheitstitel) – Einheit 1
\einheitenkopf[e1]{Einheit 1 von 4 · Punktprobe und Seitenlage}
\zweigzeile{Abitur GK}
\setcounter{aufgabe}{8}

%% TODO Titel: Ich-kann-Satz fehlt in der Bank (Platzhalter: Kettenname) – Nr. 9
%% TODO Anweisung: ein Satz über der ersten Teilaufgabe fehlt in der Bank – Nr. 9
\begin{aufgabe}{Wer und wogegen?}
\begin{teile}
\teil Gegeben sind ein Punkt $P$ und die Ebene $E: 2x + y - z = 4$. Gefragt ist, ob $P$ in $E$ liegt. Wer gegen wen? \\ \kreuz{Punkt gegen Ebene} \\ \kreuz{Gerade gegen Ebene} \\ \kreuz{Punkt gegen Gerade}
\end{teile}
\end{aufgabe}

%% TODO Titel: Ich-kann-Satz fehlt in der Bank (Platzhalter: Kettenname) – Nr. 10
%% TODO Anweisung: ein Satz über der ersten Teilaufgabe fehlt in der Bank – Nr. 10
\begin{aufgabe}{Punktprobe und Seitenlage}
%% TODO \rechenplatz{n}: Schrittzahl des Grundfalls fehlt in der Bank (nur bei fester Schreibform, 2.3 b)
\begin{teile}
\teil $E: 2x - y + z = 3$. Die Koordinaten von $P$ eingesetzt ergeben links $6$. Wo liegt $P$? \\ \kreuz{in $E$} \\ \kreuz{auf der Seite mit größerem Wert} \\ \kreuz{auf der Seite mit kleinerem Wert}
\teil Liegt $P(3 | 1 | 4)$ in $E: 2x - y + z = 7$?
\teil $E: y - 4z = -6$. Weise nach, dass $S(7 | 2 | 2)$ in $E$ liegt. \hfill (Abitur 2018 LK)
\teil $E: \vec{x} = (2 | 0 | 1) + t \cdot (1 | 1 | 0) + s \cdot (0 | 1 | 2)$. Weise nach, dass $P(4 | 5 | 7)$ in $E$ liegt. \hfill (Abitur 2021 GK)
\teil Eine Lärmschutzwand steht auf der Fahrbahn (xy-Ebene) und liegt in der Ebene $W: 5x - 2y = 0$ (1 LE = 1 m). Weise nach, dass der Punkt $K(2 | 5 | 3)$ in $W$ liegt und dass $W$ senkrecht auf der Fahrbahn steht. \hfill (Abitur 2018 GK)
\teil Ein Quader hat die Kante $AE$ mit $A(5 | 0 | 0)$ und $E(5 | 0 | 3)$; die Ebene $W: x + 2y + 2z = 9$ ist gegeben. Ein Lösungsweg lautet: (1) $P(5 | 0 | r)$ mit $0 \le r \le 3$; (2) $5 + 2 \cdot 0 + 2r = 9$. Formuliere eine passende Aufgabenstellung und führe sie zu Ende. \hfill (Abitur 2024 GK)
\teil Ein Körper ist der Teil des ersten Oktanten zwischen den parallelen Ebenen $L_1: x + 2y + z = 2$ und $L_2: x + 2y + z = 6$. Begründe, dass $T(1 | 0{,}5 | 1)$ im Inneren des Körpers liegt. \hfill (Abitur 2023 GK)
\teil Das Dreieck $ABC$ ist der Teil der Ebene $E: 2x + y + z = 8$ mit nichtnegativen Koordinaten. Aussage: Liegt $P(u | v | w)$ im Innern von $ABC$, dann auch $Q(v | u | w)$. Widerlege die Aussage mit einem Gegenbeispiel. \hfill (Abitur 2026 LK)
\end{teile}
\end{aufgabe}

%% TODO Titel: Ich-kann-Satz fehlt in der Bank (Platzhalter: Kettenname) – Nr. 11
%% TODO Anweisung: ein Satz über der ersten Teilaufgabe fehlt in der Bank – Nr. 11
\begin{aufgabe}{Punktprobe und Seitenlage – Fehler finden, Begründen}
\begin{teile}
\teil Ich finde den Fehler: $E: \vec{x} = (1 | 0 | 2) + t \cdot (1 | 2 | 0) + s \cdot (0 | 1 | 3)$, $P(3 | 5 | 8)$. Lea rechnet: aus x $t = 2$, aus z $s = 2$ – also liegt $P$ in $E$.
\teil Begründe, warum ein Punkt, dessen Koordinaten die Gleichung von $E$ erfüllen, in $E$ liegt.
\end{teile}
\end{aufgabe}
